\section{Learning and inference}
\label{sec:learning}

The model is fitted by a bidirectional variational structural EM procedure in the sense of
\citet{friedman1998learning}, specialized to a multiscale latent-variable model with exact
reconciliation. Three steps alternate at every scale. The E-step infers the latent state, the
M-step fits the parameters, and the S-step selects the lagged parent sets. One outer sweep runs
fine to coarse and then coarse to fine, for the reason
\cref{prop:model_block_invariance} gives: the block-constant component of a
scale-\(k\) predictor is identified only above \(k\) and the within-block
component only at \(k\), so each direction carries the component the other
cannot. The loop monitors one scalar objective, and rejects any sweep that
lowers it.

This section states the objective the three steps share, gives each step in the form the
implementation uses, proves the three properties of the local structure score that make the
S-step well posed, and closes with the cost of one sweep. Two of the propositions correspond to
defects that made earlier versions of this method fail, so they are stated and proved rather
than asserted.

\subsection{The penalized objective}
\label{sec:learning:objective}

Write the multiscale observations as
\[
\gD^{(k)}
=
\Big\{
\big(\vx_t^{(k)}, \vy_t^{(k)}, \gK_t^{(k)}\big)
\Big\}_{t \in \gT^{(k)}},
\qquad
\vx_t^{(k)} = \big(X_t^{(k,\ell)}\big)_{\ell \in \gL},
\quad
\vy_t^{(k)} = \big(Y_t^{(k,\ell)}\big)_{\ell \in \gL},
\]
and collect them as $\gX = \{\vx_t^{(k)}\}_{k,t}$, $\gY = \{\vy_t^{(k)}\}_{k,t}$ and
$\gK = \{\gK_t^{(k)}\}_{k,t}$. All scales are aggregations of one event log, so the
reconciliation identities hold exactly in the observed arrays. The latent variables are the
regime trajectories $\gS = \{S_t^{(k)}\}_{k,t}$ and the factor trajectories
$\gH = \{H_t^{(k)}\}_{k,t}$. The structure is the collection of lagged parent sets
$\gG = \{\gP_{k,\ell}\}$, and $\Theta$ collects the parameters of the regime chains, the factor
chains, the coarsest-scale emissions and the finer-scale allocations.

The target is the penalized marginal log-likelihood
\begin{align}
(\hat{\Theta}, \hat{\gG})
=
\argmax_{\Theta,\;\gG \in \mathfrak{G}_{\mathrm{adm}}}
\;
\log p_{\Theta,\gG}(\gX,\gY \mid \gK)
\;-\;
\sum_{k=1}^{K} \lambda_k \sum_{\ell \in \gL} \big|\gP_{k,\ell}\big|
\;-\;
\lambda_\Theta \, \Omega_\Theta(\Theta),
\label{eq:learn_penalized_marginal}
\end{align}
where $\mathfrak{G}_{\mathrm{adm}}$ is the class of admissible multiscale structures (lagged
observed-to-observed edges only, fixed top-down inter-scale edges, and same-slice co-movement
carried only by $S_t^{(k)}$ and $H_t^{(k)}$), $\lambda_k$ is the edge penalty at scale $k$, and
$\Omega_\Theta$ is a parameter regularizer.

The marginal in \cref{eq:learn_penalized_marginal} is intractable because of $\gS$ and $\gH$, so
the method optimizes the penalized evidence lower bound
\begin{align}
\gJ(q,\Theta,\gG)
=
\E_{q(\gS,\gH)}\!\big[\log p_{\Theta,\gG}(\gX,\gY,\gS,\gH \mid \gK)\big]
+
\mathbb{H}[q]
-
\sum_{k=1}^{K} \lambda_k \sum_{\ell \in \gL} \big|\gP_{k,\ell}\big|
-
\lambda_\Theta \, \Omega_\Theta(\Theta),
\label{eq:learn_elbo}
\end{align}
with $\mathbb{H}[q]$ the entropy of the variational posterior. For a fixed structure and fixed
parameters the first two terms of \cref{eq:learn_elbo} lower bound
$\log p_{\Theta,\gG}(\gX,\gY \mid \gK)$, with the gap equal to
$\KL\!\big(q \,\|\, p(\gS,\gH \mid \gX,\gY,\gK)\big)$.

The variational family factorizes top-down across scales, and within a scale it keeps the regime
chain intact and conditions the factor on it,
\begin{align}
q(\gS,\gH)
=
q^{(K)}\!\big(\gS^{(K)},\gH^{(K)}\big)
\prod_{k=1}^{K-1}
q^{(k)}\!\big(\gS^{(k)},\gH^{(k)} \mid \bar{\gZ}^{(k+1)}\big),
\qquad
q^{(k)} = q_S^{(k)}\big(\gS^{(k)}\big)\, q_H^{(k)}\big(\gH^{(k)} \mid \gS^{(k)}\big).
\label{eq:learn_variational_family}
\end{align}
The conditioning set $\bar{\gZ}^{(k+1)}$ holds what a parent scale passes to its children, namely
the parent regime marginals mapped through the parent map,
$\big(\gamma^{(k+1)}_{\parmap_k(t)}(\cdot)\big)_{t \in \gT^{(k)}}$, and the observed parent totals
$\vx_{\parmap_k(t)}^{(k+1)}$ and $\vy_{\parmap_k(t)}^{(k+1)}$. The parent factor mean
$\bar{\vh}^{(k+1)}$ is available and is not used, because a coarser factor acts on its children
through the allocation totals it has already shaped.

Three coordinate ascent steps follow from \cref{eq:learn_elbo}. Each holds two of the three
arguments fixed,
\begin{align}
q^{\new} &= \argmax_q \gJ(q,\Theta^{\old},\gG^{\old}),
\qquad
\Theta^{\new} = \argmax_\Theta \gJ(q^{\new},\Theta,\gG^{\old}),
\nonumber\\
\gG^{\new} &= \argmax_{\gG \in \mathfrak{G}_{\mathrm{adm}}} \gJ(q^{\new},\Theta^{\new},\gG).
\label{eq:learn_three_steps}
\end{align}
The E-step of \cref{eq:learn_three_steps} is solved to the accuracy of the structured family, the
M-step block by block, and the S-step label by label. \cref{sec:learning:outer} states what
survives of the monotonicity that coordinate ascent would otherwise supply.

The parameter penalty deserves one remark. In the reported fits $\lambda_\Theta$ is zero and
$\Omega_\Theta$ acts through the feasible set rather than through the objective. The factor block
solves ridged linear systems with ridge $10^{-4}$, the negative binomial dispersions are floored
at $0.05$, and the allocation concentrations are searched over $[0.5,\,2000]$. These are
constraints on where a block update may land, not an additive term, so every objective quoted
below omits $\lambda_\Theta \Omega_\Theta$.

\subsection{E-step: latent-state inference}
\label{sec:learning:estep}

Fix a scale $k$ and suppress the superscript where it is unambiguous. The scale-local objective
is
\begin{align}
\gJ_k(q_S,q_H)
=
\E_{q_S q_H}\!\big[\log p_{\Theta_k,\gG_k}\big(\gD^{(k)},\gS^{(k)},\gH^{(k)} \mid \bar{\gZ}^{(k+1)}\big)\big]
+
\mathbb{H}[q_S]
+
\E_{q_S}\!\big[\mathbb{H}[q_H(\cdot \mid \gS)]\big],
\label{eq:learn_scale_objective}
\end{align}
and it is maximized by alternating a discrete update of $q_S$ and a Gaussian update of $q_H$.

\paragraph{Expected local evidence.}
The regime update needs, for every epoch $t$ and every regime $m \in \{1,\dots,M_k\}$, the
expected log evidence
\begin{align}
\psi_t^{(k)}(m)
=
\E_{q_H}\!\left[
\log p\!\left(
\vx_t^{(k)}, \vy_t^{(k)}
\;\middle|\;
S_t^{(k)}=m,\; H_t^{(k)},\;
\big\{\gP_{k,\ell}^{X}(t),\, \gP_{k,\ell}^{Y}(t)\big\}_{\ell \in \gL}
\right)
\right].
\label{eq:learn_psi}
\end{align}
The expectation in \cref{eq:learn_psi} is evaluated at the posterior mean
$\bar{\vh}_t^{(k)} = \E_q[H_t^{(k)}]$, a first-order approximation that costs one pass over a
$T_k \times M_k \times L$ array instead of a quadrature.

At the coarsest scale the emission is a zero-inflated negative binomial count paired with a
negative binomial duration restricted to the epochs that carry at least one event,
\begin{align}
\psi_t^{(K)}(m)
&=
\sum_{\ell \in \gL}
\log \mathrm{ZINB}\!\big(X_t^{(K,\ell)} \mid \varpi_t^{(K,\ell)}(m),\, \phi_\ell^{(K)},\, \mu_t^{(K,\ell)}(m)\big)
\nonumber\\
&\quad
+
\sum_{\ell \in \gL}
\mathbf{1}\!\big\{X_t^{(K,\ell)} > 0\big\}\,
\log \mathrm{NB}\!\big(Y_t^{(K,\ell)} \mid s_\ell^{(K)},\, \nu_t^{(K,\ell)}(m)\big),
\label{eq:learn_psi_coarse}
\end{align}
with the three outputs of the emission map in
\cref{eq:model_coarse_count_mean} evaluated at
the regime indicator of $m$ and at the factor mean $\bar{\vh}_t^{(K)}$. The count mean
$\log\mu_t^{(K,\ell)}(m)$ and the duration mean $\log\nu_t^{(K,\ell)}(m)$ carry the lagged features
$\log(1+X_{t-d}^{(K,\ell')})$ and $\log(1+Y_{t-d}^{(K,\ell')})$ selected by $\gP_{K,\ell}$, the
duration mean carries the own-epoch term $\log(1+X_t^{(K,\ell)})$ as well, and the zero-inflation
predictor $\logit \varpi_t^{(K,\ell)}(m)$ carries no lagged term, so a lagged edge reaches the
count channel through the mean alone.

At a finer scale $k < K$ the emission is the allocation of a parent block. For a child epoch $t$
with parent $c = \parmap_k(t)$, let $\eta_t^{(k,\ell)}(m)$ and $\zeta_t^{(k,\ell)}(m)$ be the
count and duration allocation scores of \cref{eq:model_count_score} evaluated at
the regime indicator of $m$ and at the factor mean. Both are affine in their coefficients. The
count score carries the parent totals $\log(1+X_c^{(k+1,\ell)})$ and $\log(1+Y_c^{(k+1,\ell)})$;
the duration score carries $\log(1+Y_c^{(k+1,\ell)})$ and the own-epoch count
$\log(1+X_t^{(k,\ell)})$ in place of the parent count total. Both carry the depletion term
$\log(1+W_t^{(k,\ell)})$ represented in \cref{eq:model_count_omega}, which records what the earlier children of
the block have already taken, together with the context, the factor mean, and the lagged features
selected by $\gP_{k,\ell}$. The allocation probabilities $\omega$ and $\rho$ are block softmaxes
of these scores, and the emission is the Dirichlet-Multinomial allocation log-likelihood
\begin{align}
\psi_t^{(k)}(m)
&=
\sum_{\ell \in \gL}
\Big[
\log\Gamma\!\big(X_t^{(k,\ell)} + \kappa\,\omega_t^{(k,\ell)}(m)\big)
-
\log\Gamma\!\big(\kappa\,\omega_t^{(k,\ell)}(m)\big)
\Big]
\nonumber\\
&\quad
+
\sum_{\ell \,:\, X_t^{(k,\ell)}>0}
\Big[
\log\Gamma\!\big(Y_t^{(k,\ell)} + \tau\,\rho_t^{(k,\ell)}(m)\big)
-
\log\Gamma\!\big(\tau\,\rho_t^{(k,\ell)}(m)\big)
\Big],
\label{eq:learn_psi_fine}
\end{align}
up to a term that depends on the data and the concentrations alone and therefore shifts every
regime by the same amount. The concentrations $\kappa$ and $\tau$ carry the overdispersion of the
observed split. Letting $\kappa,\tau \to \infty$ recovers the multinomial limit
$X_t^{(k,\ell)} \log \omega_t^{(k,\ell)}$, which overstates the evidence and lets spurious lagged
edges clear the penalty.

The block softmax couples the children of one parent, so a per-regime allocation probability is
not defined without a treatment of the block-mates. The implementation scores the child under its
own regime and the block-mates at their regime-averaged score,
\begin{align}
\omega_t^{(k,\ell)}(m)
=
\frac{\exp\!\big(\eta_t^{(k,\ell)}(m)\big)}
{\exp\!\big(\eta_t^{(k,\ell)}(m)\big) + \sum_{s \in \child_k(c) \setminus \{t\}} \exp\!\big(\bar{\eta}_s^{(k,\ell)}\big)},
\qquad
\bar{\eta}_s^{(k,\ell)} = \sum_{m'=1}^{M_k} \gamma_s(m')\, \eta_s^{(k,\ell)}(m'),
\label{eq:learn_loo_omega}
\end{align}
and likewise for $\rho$ with the score $\zeta$ and the active-child restriction. This mean-field
treatment of the block normalizer is the one place where the objective a step is accepted on
differs from the objective the next step reports, and \cref{sec:learning:outer} handles the
consequence.

\paragraph{Regime update.}
With $\psi^{(k)}$ in hand, $q_S^{(k)}$ is the chain distribution
\begin{align}
q_S\big(\gS^{(k)}\big)
\;\propto\;
p\big(S_1^{(k)} \mid \gK_1^{(k)}\big)
\prod_{t=2}^{T_k}
p\big(S_t^{(k)} \mid S_{t-1}^{(k)}, \gK_t^{(k)}, \bar{\gZ}^{(k+1)}\big)
\prod_{t=1}^{T_k}
\exp\!\big(\psi_t^{(k)}(S_t^{(k)})\big),
\label{eq:learn_qS}
\end{align}
whose transitions replace the parent regime of \cref{eq:model_regime_init} by its posterior
marginal,
\begin{align}
p\big(S_t^{(k)} = m \mid S_{t-1}^{(k)} = m',\, \bar{\gZ}^{(k+1)}\big)
=
\softmax_m\!\Big(
a_m^{(k)} + A_{m,m'}^{(k)}
+ \sum_{n=1}^{M_{k+1}} \gamma^{(k+1)}_{\parmap_k(t)}(n)\, C_{m,n}^{(k)}
+ (\vb_m^{(k)})^\top \gK_t^{(k)}
\Big).
\label{eq:learn_trans_logit}
\end{align}
The cross-scale term enters through the expected logit. That form is exact when the parent
posterior is degenerate, and it is the same first-order form the M-step differentiates, so the
two steps optimize one objective. Because \cref{eq:learn_qS} is a Markov chain with time-varying
transitions, a log-space forward-backward recursion returns the marginals $\gamma_t(m)$, the pair
marginals $\xi_t(m,m')$ and the log evidence $\log Z_k$ exactly, at cost $O(T_k M_k^2)$.

\paragraph{Factor update.}
Conditioned on the regime path the factor dynamics are linear Gaussian, but the emissions are
not, so $q_H$ is taken Gaussian and fitted by a regime-averaged Kalman filter followed by an RTS
smoother. The averaged system at epoch $t$ is
\begin{align}
\bar{\mF}_t &= \sum_{m} \gamma_t(m)\, \mF_m,
\qquad
\bar{\vg}_t = \sum_{m} \gamma_t(m)\, \mG_m \gK_t^{(k)},
\nonumber\\
\bar{\mQ}_t &= \sum_{m} \gamma_t(m)\Big[\mSigma_m + (\mF_m - \bar{\mF}_t)\, \mP_{t-1} (\mF_m - \bar{\mF}_t)^\top\Big],
\label{eq:learn_avg_lgssm}
\end{align}
where $\mP_{t-1}$ is the filtered covariance of the previous epoch. The second term of
$\bar{\mQ}_t$ is the Jensen gap that a single averaged transition matrix would otherwise discard.

At the coarsest scale the emission carries information about the factor, and that information is
injected as a Laplace pseudo-observation. Write $\psi_t^{(K)}(m;\vh)$ for
\cref{eq:learn_psi_coarse} with the factor held at an arbitrary $\vh$ rather than at
$\bar{\vh}_t^{(K)}$, and let
\[
\bar{\psi}_t(\vh)
=
\sum_{m} \gamma_t(m)\, \psi_t^{(K)}(m;\vh)
\]
be the regime-averaged emission log-likelihood at epoch $t$ as a function of the factor. Its
gradient and Hessian at the current mean are taken by central differences, the Hessian eigenvalues
are clipped above at $-10^{-6}$ to force negative definiteness, and the filter absorbs
\begin{align}
\mR_t = \big(-\nabla^2 \bar{\psi}_t(\bar{\vh}_t)\big)^{-1},
\qquad
\vz_t = \bar{\vh}_t + \mR_t\, \nabla \bar{\psi}_t(\bar{\vh}_t)
\label{eq:learn_laplace}
\end{align}
as an observation of $H_t^{(K)}$ under the identity observation model with noise covariance
$\mR_t$. This is one Newton step on the emission, converted into the moments a Kalman update
consumes. At a finer scale no pseudo-observation is formed and the smoother runs on the prior
alone, since the factor there enters only the allocation scores, whose curvature is already
represented through the block normalizer.

\paragraph{Alternation.}
The two updates are alternated. The regime update needs $\bar{\vh}$ to build $\psi$, and the
factor update needs $\gamma$ to build $\bar{\mF}_t$ and $\bar{\psi}_t$. The loop starts from
$\bar{\vh} = \vzero$ and a uniform $\gamma$, and stops when
$|\log Z_k - \log Z_k^{\mathrm{prev}}| < \varepsilon\,(1 + |\log Z_k|)$ or after a fixed
iteration budget. One case is exact in a single pass. At the coarsest scale with $r_K = 0$
nothing feeds back into the emissions, so the loop exits immediately. At a finer scale with
$r_k = 0$ the loop still runs, because the block normalizer of \cref{eq:learn_loo_omega} depends
on $\gamma$ through $\bar{\eta}$, so the two updates still have a fixed point to reach.

\subsection{M-step: parameter learning}
\label{sec:learning:mstep}

Given $q^{\new}$ and the current structure, the parameters solve
\begin{align}
\Theta_k^{\new}
=
\argmax_{\Theta_k}
\;
\E_q\!\big[\log p_{\Theta_k,\gG_k}\big(\gD^{(k)},\gS^{(k)},\gH^{(k)} \mid \bar{\gZ}^{(k+1)}\big)\big].
\label{eq:learn_mstep}
\end{align}
Given $q$, the expected complete-data log-likelihood separates into parameter blocks that share no
arguments, so the blocks can be updated in any order. Five blocks appear in total, and one scale
uses the regime block, the factor block, and then either the two coarsest-scale emission blocks or
the single allocation block.

\begin{enumerate}
\item \textbf{Regime transitions} $\{\va^{(k)}, A^{(k)}, C^{(k)}, \vb^{(k)}\}$ together with the
initial distribution. The objective is the expected log-probability of the regime path under
\cref{eq:learn_trans_logit}, a softmax conditional likelihood whose gradient is the difference
between expected and predicted transition mass,
\[
\frac{\partial}{\partial\, \logit_t(m \mid m')}
=
\xi_{t-1}(m',m) - \gamma_{t-1}(m')\, p_t(m \mid m').
\]
It is maximized by Adam. A tabular soft-Bayes estimate of the same block can be enabled as a
second proposal, in which case the proposal with the higher block objective is taken; the fits
reported here use the gradient proposal alone.

\item \textbf{Latent factor} $\{\mF_m, \mG_m, \mSigma_m, \vmu_0, \mSigma_0\}$. The expected
complete-data log-likelihood of a switching linear Gaussian chain is quadratic in $\mF_m$ and
$\mG_m$ given the smoothed moments, so the update is closed form, a regime-weighted least squares
solve per regime with ridge $10^{-4}$, followed by the weighted residual second moment for
$\mSigma_m$. One approximation is made here. The cross moment $\E_q[H_t H_{t-1}^\top]$ is replaced
by $\bar{\vh}_t \bar{\vh}_{t-1}^\top$, dropping the smoother cross-covariance, so the update is
exact in the means and optimistic in the covariances.

\item \textbf{Coarsest-scale counts.} The zero-inflated negative binomial coefficients and the
per-label dispersion are fitted by Adam on the $\gamma$-weighted log-likelihood of
\cref{eq:learn_psi_coarse}, with a zero observation contributing to the mean gradient through its
posterior responsibility for coming from the negative binomial rather than the inflation
component.

\item \textbf{Coarsest-scale durations.} The negative binomial duration coefficients and
dispersions are fitted by Adam on the same weighted objective, restricted to the epochs with a
positive count.

\item \textbf{Finer-scale allocations.} The count and duration score coefficients are fitted by
Adam on the expected Dirichlet-Multinomial allocation log-likelihood, whose gradient with respect
to a score is the block-centered digamma residual
\[
\frac{\partial}{\partial \eta_t^{(k,\ell)}}
=
\kappa\, \omega_t^{(k,\ell)} \Big[ v_t - \sum_{s \in \child_k(c)} \omega_s^{(k,\ell)} v_s \Big],
\qquad
v_t = \Psi\!\big(X_t^{(k,\ell)} + \kappa \omega_t^{(k,\ell)}\big) - \Psi\!\big(\kappa \omega_t^{(k,\ell)}\big),
\]
with $\Psi$ the digamma function. The concentrations $\kappa$ and $\tau$ are estimated afterwards,
one per channel, by golden-section search on the log concentration over $[0.5,\,2000]$, using the
complete Dirichlet-Multinomial
log-likelihood including the block-constant term that the per-epoch form of
\cref{eq:learn_psi_fine} drops, since that term is the only part of the likelihood that
identifies a concentration. Leaving $\kappa$ and $\tau$ at an assumed value makes every structure
score that uses this likelihood over- or under-confident, so they are fitted rather than fixed.
\end{enumerate}

A fixed number of Adam steps can overshoot, so each gradient-based update is accepted only if it
does not lower the expected log-likelihood it targets, and otherwise the previous value is kept.
The regime block is checked on its own; the two coarsest-scale emission blocks are proposed
together and accepted or rejected together, and so are the two allocation blocks. The factor block
is a closed-form solve rather than a gradient step and is applied without a check. Each checked
group is therefore monotone in its own objective. That does not make the sweep monotone in the
objective the outer loop monitors, for the reason given in \cref{sec:learning:outer}.

\subsection{S-step: structure learning}
\label{sec:learning:sstep}

For each scale $k$ and label $\ell$, let
$\gC_{k,\ell} \subseteq \gL \times \{1,\dots,p_k\}$ be the candidate set of lagged parents, where
$p_k$ is the maximum lag in \cref{eq:model_lag_values}. With no screening rule the candidate
set is the whole product, so $|\gC_{k,\ell}| = L\,p_k$. Write $D_k$ for the maximum in-degree
allowed at scale $k$, a quantity distinct from the context dimension $d_k$. The learned parent set
satisfies $\gP_{k,\ell} \subseteq \gC_{k,\ell}$ and $|\gP_{k,\ell}| \le D_k$.

\paragraph{The local score.}
Let $\vbeta_\gP$ collect the lag coefficients that a parent set $\gP$ contributes to the two linear
predictors of label $\ell$, namely $\beta_{\ell',d}^{(k,\ell)}$ in the count channel and
$\alpha_{\ell',d}^{(k,\ell)}$ in the duration channel for each $(\ell',d) \in \gP$, and let
$\Theta_{-\gP}$ be everything else, held at its current value. The local score is the expected
local log-likelihood \emph{maximized over} $\vbeta_\gP$,
\begin{align}
Q_{k,\ell}(\gP; q,\Theta)
=
\max_{\vbeta_\gP}
\;
\sum_{t \in \gT^{(k)}} \sum_{m=1}^{M_k}
\gamma_t(m)\,
\log p_{\vbeta_\gP,\Theta_{-\gP}}\!\left(
X_t^{(k,\ell)}, Y_t^{(k,\ell)}
\;\middle|\;
S_t^{(k)}=m,\; \bar{\vh}_t^{(k)},\; \gP^{X}(t),\, \gP^{Y}(t)
\right),
\label{eq:learn_local_score}
\end{align}
with $\gP^{X}(t)$ and $\gP^{Y}(t)$ the lagged parent values of \cref{eq:model_lag_values}, and the
structure update at $(k,\ell)$ is
\begin{align}
\hat{\gP}_{k,\ell}
=
\argmax_{\gP \subseteq \gC_{k,\ell},\; |\gP| \le D_k}
\;
Q_{k,\ell}(\gP; q,\Theta) - \lambda_k |\gP|.
\label{eq:learn_local_update}
\end{align}
At the coarsest scale the density in \cref{eq:learn_local_score} is the zero-inflated negative
binomial paired with the negative binomial duration, and the implementation evaluates the
$\gamma$-weighted sum exactly as written. At a finer scale it is the Dirichlet-Multinomial
allocation likelihood of \cref{eq:learn_psi_fine} with the block normalizer included, and the
implementation makes one further economy. Rather than form the leave-one-out allocation of
\cref{eq:learn_loo_omega} once per regime, it evaluates the allocation likelihood once at the
regime-averaged score $\bar{\eta}$, which is the $\gamma$-weighted sum of
\cref{eq:learn_local_score} to first order and coincides with it when the regime posterior is
degenerate. The gradient that drives the inner maximization is still contracted against $\gamma$.
Two features of \cref{eq:learn_local_score} are load-bearing, the block normalizer and the
maximization over $\vbeta_\gP$. Both were absent from earlier versions of the method, and the next
two propositions say exactly what their absence costs.

\begin{proposition}[An unnormalized local score has no maximizer]
\label{prop:learn_unnormalized}
Fix a scale $k < K$, a label $\ell$ and a parent set $\gP$. Let $\vz_t \in \sR^{J}$ collect the
features entering the count allocation score of epoch $t$, so that
$\eta_t^{(k,\ell)} = \vtheta^\top \vz_t$ for a coefficient vector $\vtheta \in \sR^{J}$, and
define the unnormalized local score
\begin{align}
\widetilde{Q}_{k,\ell}(\gP;\vtheta)
=
\sum_{t \in \gT^{(k)}} X_t^{(k,\ell)}\, \eta_t^{(k,\ell)}
=
\vtheta^\top \vw,
\qquad
\vw = \sum_{t \in \gT^{(k)}} X_t^{(k,\ell)}\, \vz_t .
\label{eq:learn_unnormalized_score}
\end{align}
If $\vw \neq \vzero$ then
$\sup_{\vtheta} \big[\widetilde{Q}_{k,\ell}(\gP;\vtheta) - \lambda_k |\gP|\big] = +\infty$ and the
supremum is not attained. Consequently, for any two parent sets whose feature sums are both
nonzero, the profiled unnormalized scores are both infinite, and their comparison carries no
information about the parent sets.
\end{proposition}

\begin{proof}
$\widetilde{Q}_{k,\ell}$ is linear in $\vtheta$, so for $c>0$ it satisfies
$\widetilde{Q}_{k,\ell}(\gP; c\vtheta) = c\,\vtheta^\top \vw$. Take $\vtheta = \vw$; then
$\widetilde{Q}_{k,\ell}(\gP; c\vw) = c\|\vw\|^2 \to \infty$ as $c \to \infty$, because
$\|\vw\|^2 > 0$. The penalty $-\lambda_k|\gP|$ is a finite constant and does not affect the
divergence, and no finite $\vtheta$ attains an infinite value. The same argument applies to any
other parent set with a nonzero feature sum, so \cref{eq:learn_local_update} would compare
$+\infty$ with $+\infty$.
\end{proof}

\begin{remark}
\label{rem:learn_normalizer_bounds}
The block-normalized score is bounded above and the difficulty disappears. Under the multinomial
allocation, $\sum_t X_t^{(k,\ell)} \log \omega_t^{(k,\ell)} \le 0$ because each
$\omega_t^{(k,\ell)} \in (0,1]$. Under the Dirichlet-Multinomial form of
\cref{eq:learn_psi_fine}, the per-epoch term satisfies, for integer $X \ge 0$ and
$a = \kappa\omega \in (0,\kappa]$,
\[
\log\Gamma(X+a) - \log\Gamma(a) = \sum_{j=0}^{X-1} \log(a+j) \;\le\; \sum_{j=0}^{X-1} \log(\kappa+j),
\]
since the middle expression is nondecreasing in $a$ and the right side is its value at
$\omega = 1$. The right side depends on the data and on $\kappa$ alone, so summing over epochs and
labels, and repeating the argument on the duration channel with $\tau$, bounds the whole score by
a constant free of the allocation coefficients. The supremum of
\cref{eq:learn_local_score} is therefore finite and two parent sets are compared as two finite
numbers. The block normalizer is what supplies the bound, because it is the only term that
charges for raising one child's score without lowering its block-mates.
\end{remark}

\begin{proposition}[Scoring at the incumbent coefficient can only delete edges]
\label{prop:learn_profiling}
Let $\widehat{Q}_{k,\ell}(\gP;\Theta)$ denote the local expected log-likelihood of
\cref{eq:learn_local_score} evaluated at the current $\Theta$ rather than maximized over
$\vbeta_\gP$. Fix $\gP \subseteq \gC_{k,\ell}$ and suppose the parameterization holds the count and
duration lag coefficients of every candidate outside $\gP$ at zero, that is
$\beta_e^{(k,\ell)} = \alpha_e^{(k,\ell)} = 0$ for every $e \in \gC_{k,\ell} \setminus \gP$, which
is what a parameterization indexed by the current parent set supplies when $\gP$ is that set. Let
$\lambda_k > 0$.
\begin{enumerate}
\item For every $e \in \gC_{k,\ell} \setminus \gP$,
\begin{align}
\Big[\widehat{Q}_{k,\ell}(\gP \cup \{e\};\Theta) - \lambda_k|\gP \cup \{e\}|\Big]
-
\Big[\widehat{Q}_{k,\ell}(\gP;\Theta) - \lambda_k|\gP|\Big]
=
-\lambda_k \;<\; 0 .
\label{eq:learn_addition_is_minus_lambda}
\end{align}
A search that accepts a move only when the penalized score strictly increases therefore never
adds an edge, and the parent set it returns is a subset of the one it started from.
\item The profiled score $Q_{k,\ell}$ of \cref{eq:learn_local_score} is monotone under enlargement
of the parent set,
\begin{align}
Q_{k,\ell}(\gP \cup \{e\}) \;\ge\; Q_{k,\ell}(\gP),
\qquad\text{so}\qquad
Q_{k,\ell}(\gP \cup \{e\}) - Q_{k,\ell}(\gP) - \lambda_k \;\ge\; -\lambda_k ,
\label{eq:learn_profiled_gain}
\end{align}
with equality exactly when the extra coefficient buys no improvement in fit, and strict positivity
exactly when the best achievable improvement exceeds $\lambda_k$.
\end{enumerate}
\end{proposition}

\begin{proof}
For part 1, the parent set enters the local log-likelihood only through the lagged term
$\sum_{(\ell',d) \in \gP} \beta_{\ell',d}^{(k,\ell)} \log\!\big(1 + X_{t-d}^{(k,\ell')}\big)$ of
the count predictor and its counterpart
$\sum_{(\ell',d) \in \gP} \alpha_{\ell',d}^{(k,\ell)} \log\!\big(1 + Y_{t-d}^{(k,\ell')}\big)$ in
the duration predictor. Adding $e = (\ell',d)$ with $\beta_e^{(k,\ell)} = \alpha_e^{(k,\ell)} = 0$
adds the terms $0 \cdot \log(1+X_{t-d}^{(k,\ell')}) = 0$ and $0 \cdot \log(1+Y_{t-d}^{(k,\ell')})
= 0$ to the predictors of every epoch. The predictors are therefore unchanged, hence so is every
emission factor, hence so is $\widehat{Q}_{k,\ell}$. The penalty grows by exactly $\lambda_k$,
which gives \cref{eq:learn_addition_is_minus_lambda}. A hill climb that keeps the parameterization
indexed by its current parent set meets the hypothesis at every state it reaches, so no addition
ever increases the penalized score, the only accepted moves are removals, and the returned set is
contained in the initial set.

For part 2, $\vbeta_\gP$ ranges over $\sR^{2|\gP|}$, two coefficients per edge. The map
$\vbeta_\gP \mapsto (\vbeta_\gP, \vzero)$ embeds $\sR^{2|\gP|}$ into $\sR^{2|\gP|+2}$, the
coefficient space of $\gP \cup \{e\}$, and leaves the objective value unchanged, by the same
argument as part 1. A maximum over a superset is at least the maximum over the subset, which gives
\cref{eq:learn_profiled_gain}. Equality holds exactly when the maximizer over the larger space
attains no more than the maximizer over the smaller one, that is, when the additional coefficients
buy nothing.
\end{proof}

\begin{remark}
\label{rem:learn_profile_practice}
Profiling the candidate's coefficients before comparing is therefore what makes an addition
possible at all, and it is what gives $\lambda_k$ its intended meaning as the improvement in fit
that a new coefficient must buy. The implementation approximates the inner maximum of
\cref{eq:learn_local_score} by normalized-gradient ascent from the incumbent coefficients, with a
backtracking line search that accepts only improving steps, so the returned value is never below
the value at the starting point. The monotonicity of \cref{eq:learn_profiled_gain} therefore
holds exactly for the true maximizer and up to the accuracy of the ascent in practice. The fitted
coefficients of the winning parent set are written back into $\Theta$, so a newly accepted edge
starts from a fitted value rather than from zero.
\end{remark}

\begin{proposition}[Per-label decomposition]
\label{prop:learn_decomposition}
Fix $q$ and $\Theta$.
\begin{enumerate}
\item For every assignment $(\gP_{k,\ell}) \in \prod_{k,\ell} 2^{\gC_{k,\ell}}$, the unrolled
multiscale graph is acyclic.
\item Write $\gJ^{\ast}(q,\Theta,\gG)$ for the value of $\gJ$ after maximization over the lag
coefficients $\{\vbeta_{\gP_{k,\ell}}\}_{k,\ell}$ alone, with the rest of $\Theta$ held fixed.
Then the dependence of $\gJ^{\ast}$ on $\gG$ is a sum of per-label terms,
\begin{align}
\gJ^{\ast}(q,\Theta,\gG)
=
R(q,\Theta)
+
\sum_{k=1}^{K} \sum_{\ell \in \gL}
\Big[ Q_{k,\ell}(\gP_{k,\ell}; q,\Theta) - \lambda_k |\gP_{k,\ell}| \Big],
\label{eq:learn_separable}
\end{align}
where $R$ does not depend on $\gG$.
\item Consequently the joint structure problem splits into $K L$ independent subset-selection
problems, and the collection of per-label maximizers of \cref{eq:learn_local_update} is a
maximizer of $\gJ^{\ast}$ over $\mathfrak{G}_{\mathrm{adm}}$.
\end{enumerate}
\end{proposition}

\begin{proof}
For part 1, give each node of the unrolled graph the key $\big(-k,\; t,\; \sigma\big)$, compared
lexicographically, where $k$ is the scale, $t$ is the epoch index within that scale, and
$\sigma \in \{0,1,2,3\}$ takes the value $0$ on the regime node $S_t^{(k)}$, $1$ on the factor node
$H_t^{(k)}$, $2$ on the count nodes $X_t^{(k,\ell)}$, and $3$ on the duration nodes
$Y_t^{(k,\ell)}$. Every admissible edge strictly increases this key. An inter-scale edge runs from
scale $k+1$ to scale $k$ and increases the first component, and no edge runs the other way. A
within-scale edge between two epochs runs from epoch $t-d$ to epoch $t$ with $d \ge 1$, leaving the
first component fixed and increasing the second, which covers the lagged observed edges, the two
latent chains, and the depletion feature, since the latter reads only children of the same block
that precede $t$. Within one epoch, $S_t^{(k)} \to H_t^{(k)}$ raises $\sigma$ from $0$ to $1$, a
latent-to-observed edge raises $\sigma$ from $0$ or $1$ to $2$ or $3$, and the count-to-duration
edge $X_t^{(k,\ell)} \to Y_t^{(k,\ell)}$ raises $\sigma$ from $2$ to $3$. A directed cycle would
have to return to its starting key, which a strictly increasing sequence cannot do. Since the
candidate sets contain no lag-zero pair, the argument covers every assignment, and no acyclicity
constraint has to be imposed on the search.

For part 2, the expected complete-data log-likelihood at scale $k$ splits into a regime-chain
term, a factor-chain term, and the per-epoch emission terms of \cref{eq:learn_psi_coarse} or
\cref{eq:learn_psi_fine}. The first two do not involve $\gG$. The emission terms are sums over
labels, and $\gP_{k,\ell}$ enters only the linear predictor of label $\ell$. At a finer scale the
block normalizer of \cref{eq:learn_loo_omega} normalizes over the children of one parent block
for one fixed label, so it couples epochs within a label but never two labels. The entropy
$\mathbb{H}[q]$ and the parameter penalty do not involve $\gG$. The coefficients
$\vbeta_{\gP_{k,\ell}}$ appear in the label-$\ell$ terms of scale $k$ and in no other term, so the
maximization defining $\gJ^{\ast}$ distributes over the labels and scales, and collecting the
label-$\ell$ terms and adding the edge penalty gives \cref{eq:learn_separable} with the profiled
local score of \cref{eq:learn_local_score}.

For part 3, the feasible set
$\prod_{k,\ell} \{\gP \subseteq \gC_{k,\ell} : |\gP| \le D_k\}$ is a product and, by part 1,
carries no joint constraint. A sum of terms over a product set is maximized coordinatewise, so
maximizing each bracket of \cref{eq:learn_separable} separately maximizes the whole.
\end{proof}

\begin{corollary}[Exact per-target search is feasible]
\label{cor:learn_exact_search}
Under the in-degree bound $D_k$, an exhaustive search of one subproblem in
\cref{eq:learn_local_update} evaluates
$\sum_{j=0}^{D_k} \binom{|\gC_{k,\ell}|}{j} = O\!\big(|\gC_{k,\ell}|^{D_k}\big)$ parent sets, which
is polynomial in the candidate-set size for fixed $D_k$. By \cref{prop:learn_decomposition},
solving the $L$ subproblems at a scale exactly solves the structure problem at that scale exactly,
with no acyclicity projection and no search over variable orderings.
\end{corollary}

The decomposability in \cref{eq:learn_separable} is the property that score-based Bayesian network
learning relies on generally \citep{heckerman1995learning,chickering2002optimal}; what is added
here is that the lagged restriction removes the acyclicity constraint that normally couples the
per-target problems back together \citep{murphy2002dbn,pamfil2020dynotears}.

\paragraph{Search.}
The default search is greedy forward-backward hill climbing, in the manner of the local search
stage of max-min hill climbing \citep{tsamardinos2006max}. Each pass evaluates every single-edge
addition that respects the in-degree bound and every single-edge removal, accepts the moves that
improve the penalized score, and stops when a pass finds no improvement. Scores are memoized on
the parent set, since the passes revisit the same sets.

Greedy hill climbing accepts one edge at a time, so a pair of edges that pays only jointly is
unreachable from the incumbent set. Three alternatives are implemented against that. Beam search
carries a fixed-width frontier of parent sets of increasing size. Exhaustive enumeration is exact,
and by \cref{cor:learn_exact_search} it is affordable whenever the subset count within the
in-degree bound is small, which holds for every fit reported here. A genetic search treats a
parent set as an individual and recombines and mutates them, for candidate sets too large to
enumerate. \Cref{sec:results:search} reports that all three return the same parent sets as the
greedy search on both benchmarks, so on these data the greedy answer is already the certified
per-target optimum.

The edge penalty is $\lambda_k = \lambda \log T_k$, the Bayesian information criterion price of
the two coefficients (one count, one duration) that an edge adds at a scale with $T_k$ epochs, so
that a single $\lambda$ behaves consistently across scales and sample sizes. The value of
$\lambda$ is not a constant of the method. It is selected together with the latent orders on
held-out log evidence, as described next.

\paragraph{Selecting the latent orders and the penalty.}
The regime count $M_k$, the factor dimension $r_k$ and the penalty multiplier $\lambda$ are free
quantities of the model rather than quantities a single fit determines. Each candidate triple is
fitted on a training split by the identical procedure and scored by
$-\sum_k \log Z_k$ on a held-out split, and the triple with the lowest score is kept. Nothing else
is selected this way. The graph budget $(\text{max lag}, D_k)$ and the block lengths of the
hierarchy state which lags and which timescales are of interest, and a held-out score cannot
supply a question, so both remain declarations that the protocol discloses.

\subsection{The outer loop and its monotonicity guard}
\label{sec:learning:outer}

One sweep is two passes. The fine-to-coarse pass runs the three steps at every
scale from $k=1$ up to $k=K$; the coarse-to-fine pass runs them from $k=K$ back
down to $k=1$, passing $\bar{\gZ}^{(k+1)}$ to each child scale.
\Cref{prop:model_block_invariance} is what makes both necessary. A scale-$k$
predictor splits into a block-constant part, which the child kernel is blind to
and which scale $k+1$ identifies, and a within-block part, which scale $k$
identifies. Each pass carries the part the other leaves unidentified. The
monitored objective is
\begin{align}
\gJ^{\mathrm{sweep}}
=
\sum_{k=1}^{K} \log Z_k
-
\sum_{k=1}^{K} \lambda_k \sum_{\ell \in \gL} \big|\gP_{k,\ell}\big|,
\label{eq:learn_monitored}
\end{align}
the summed per-scale log evidence returned by the forward-backward recursion, less the edge
penalty the structure step trades against it, under the same penalty schedule
$\lambda_k = \lambda \log T_k$ the structure step uses. The log evidence $\log Z_k$ is the
scale-local bound on the regime chain at the current expected local evidence, so it already
carries the entropy of $q_S^{(k)}$. Tracking the log evidence alone would leave out the
term the structure step trades against, and a sweep that correctly drops an edge lowers
$\sum_k \log Z_k$ while raising $\gJ^{\mathrm{sweep}}$, so the log evidence on its own would read
as a step backwards. The convergence
test is relative. The fit stops once the last two changes in \cref{eq:learn_monitored} each
satisfy $|\Delta \gJ^{\mathrm{sweep}}| < \varepsilon\,(1 + |\gJ^{\mathrm{sweep}}|)$. An absolute
test cannot fire when the objective is in the hundreds or thousands, and a fit under an absolute
test runs its full sweep budget however long it has been converged.

Coordinate ascent on \cref{eq:learn_elbo} would make \cref{eq:learn_monitored} nondecreasing, and
it does not, for one identifiable reason. The allocation normalizer of a finer scale is evaluated
at the regime-averaged score, as in \cref{eq:learn_loo_omega}. The M-step accepts an allocation
block when the allocation log-likelihood computed at
$\bar{\eta} = \sum_m \gamma(m)\,\eta(m)$ does not decrease. The next E-step reports $\log Z_k$
from a likelihood in which the child under evaluation carries its own regime's score while its
block-mates carry $\bar{\eta}$. The two quantities agree when the regime posterior is degenerate
and differ otherwise, so the quantity a parameter block is accepted on is not the quantity the
next inference step reports, and a sweep in which every block passed its own check can still lower
\cref{eq:learn_monitored}.

The outer loop turns that residual into a guard rather than a decreasing trace. A sweep whose
objective falls below the best sweep so far is rejected. The state (parameters, graph and
posterior summaries) reverts to the best sweep and the sweep is retried with the M-step learning
rate halved. Two consecutive rejections are retried this way and a third stops the fit at its best
state; an accepted sweep resets the count. A rejection whose
shortfall is itself below the convergence tolerance is read as convergence rather than as failure
and stops the fit immediately. The fit returns its best sweep rather than its last one, and the
rejected objectives are recorded separately, so the rejections stay measurable while the recorded
trace stays monotone. The five simulation experiments reported in this monograph record 1176 runs,
of which 1056 drive the outer loop; the remaining 120 are the oracle-parameter reference, which
runs one E-step and no sweep. Of those 1056, 860 triggered the guard at least once, and 1893
sweeps were rejected in total. The residual the guard absorbs is therefore a routine feature of
the fit rather than an edge case.

\subsection{Shift confirmation of accepted edges}
\label{sec:learning:confirm}

The penalty $\lambda_k$ is a fixed price, and at a scale where a persistent latent regime or a
shared parent-block total autocorrelates every label, a lagged feature that carries no information
about the target can still clear it. Each edge the search accepts is therefore put through a
second test against a null built from its own feature.

Fix an accepted edge $e \in \hat{\gP}_{k,\ell}$ and let
$\gP_{-e} = \hat{\gP}_{k,\ell} \setminus \{e\}$. The realized gain of the edge is
$G_e = Q_{k,\ell}(\hat{\gP}_{k,\ell}) - Q_{k,\ell}(\gP_{-e})$, penalties excluded. For
$b = 1,\dots,B$, a surrogate $\tilde{e}^{(b)}$ is built by displacing the rows of the design
column of $e$, and its gain over the same base set,
$G^{(b)} = Q_{k,\ell}\big(\gP_{-e} \cup \{\tilde{e}^{(b)}\}\big) - Q_{k,\ell}(\gP_{-e})$, is
recorded. The edge is kept only if $G_e$ exceeds an upper quantile of $\{G^{(b)}\}_{b=1}^{B}$ at a
level fixed in advance. The penalty cancels from both sides of that comparison, since the real
edge and every surrogate enlarge the base set by one. The experiments reported here use $B = 32$
copies and the $95\%$ quantile.

The displacement is chosen so that the null retains everything about the feature except its
alignment with the target. At the coarsest scale it is a whole-series circular shift by a random
offset drawn away from the endpoints, which preserves the marginal distribution of the feature and
its autocorrelation. A plain row shuffle would destroy the autocorrelation as well and would
understate the floor, because a lagged feature is attractive to the score precisely for being
autocorrelated, which is also what the latent regime explains. At a finer scale the column is
rolled independently inside each parent block, so block membership survives into the null. Gains
that come from the shared block total, from a persistent latent state, or from the realized
allocation split of a block therefore appear in the null too, and only epoch-level alignment can
separate a real edge from it. The test runs to a fixed point, since removing one edge changes the
base set of the survivors and hence their gains, so the survivors are rechecked after each
removal.

The scope of the test is worth stating plainly. It removes edges whose apparent gain is matched by
a feature with the same marginal distribution, the same autocorrelation and, at a finer scale, the
same block membership. It does not remove an edge whose source is genuinely aligned with the
target at the epoch level, whether or not that alignment is causal, so a common driver that is
itself epoch-aligned survives. It applies only to the edges the search proposed, so it can lower
recall but can never raise it, and it makes no correction for the number of candidates the search
examined before proposing them. It costs $B$ additional score evaluations per accepted edge per
pass rather than per candidate.

\subsection{Computational and memory complexity}
\label{sec:learning:complexity}

Write $T_k$ for the number of epochs at scale $k$, $L = |\gL|$ for the number of labels, $M_k$ for
the number of regimes, $r_k$ for the factor dimension, $d_k$ for the context dimension, $D_k$ for
the in-degree bound, and $\bar{C}_k = L^{-1}\sum_{\ell \in \gL} |\gC_{k,\ell}|$ for the average
candidate-set size, which equals $L p_k$ when nothing is screened. Let $I_E$ be the number of
E-step alternations, $n_\Theta$ the number of M-step gradient steps per block, $n_Q$ the number of
gradient steps used to profile one candidate parent set, and $B$ the number of shifted copies per
accepted edge.

\paragraph{E-step.}
The lagged design contributions are built once per E-step, at cost $O(T_k L \bar{C}_k)$. Each
alternation then evaluates the emission arrays, which have shape $T_k \times M_k \times L$, runs
the forward-backward recursion, and runs the smoother,
\begin{align}
O\!\Big(
T_k L \bar{C}_k
\;+\;
I_E \big[\, T_k M_k L + T_k M_k^2 + T_k r_k^3 \,\big]
\Big).
\label{eq:learn_cost_estep}
\end{align}
At the coarsest scale with $r_K > 0$ the Laplace step of \cref{eq:learn_laplace} adds a finite
difference gradient and Hessian, which take $2 r_K$ and $2 r_K (r_K + 1)$ evaluations of
$\bar{\psi}$ respectively, plus one eigendecomposition per epoch,
\begin{align}
O\!\Big(
I_E \big[\, T_K r_K^2 M_K L + T_K r_K^3 \,\big]
\Big).
\label{eq:learn_cost_laplace}
\end{align}

\paragraph{M-step.}
The gradient-based blocks cost one pass over the same arrays per step, and the factor block is a
closed-form solve per regime,
\begin{align}
O\!\Big(
n_\Theta \big[\, T_k M_k L + T_k L \bar{C}_k + T_k M_k^2 \,\big]
\;+\;
M_k r_k^3 + T_k M_k r_k^2
\;+\;
T_k L
\Big),
\label{eq:learn_cost_mstep}
\end{align}
where the last term is the golden-section search for the concentrations, whose iteration count is
fixed.

\paragraph{S-step.}
Evaluating \cref{eq:learn_local_score} for one parent set of size at most $D_k$ profiles at most
$D_k$ coefficients per channel, and each inner step touches a $T_k \times M_k$ array and a
$T_k \times D_k$ design matrix, so one score evaluation costs $O\!\big(n_Q T_k (M_k + D_k)\big)$.
Greedy hill climbing evaluates $O(D_k |\gC_{k,\ell}|)$ parent sets per label, since the in-degree
bound caps the number of productive passes. Exhaustive enumeration under the same bound evaluates
$O\!\big(|\gC_{k,\ell}|^{D_k}\big)$ of them. Shift confirmation adds $O(B D_k^2)$ evaluations per
label, since each of at most $D_k$ accepted edges draws $B$ surrogates and every removal restarts
the check. The S-step at scale $k$ therefore costs
\begin{align}
O\!\Big(
n_Q T_k (M_k + D_k)\; L \big[\, D_k \bar{C}_k + B D_k^2 \,\big]
\Big)
\quad\text{greedy},
\qquad
O\!\Big(
n_Q T_k (M_k + D_k)\; L \, \bar{C}_k^{\,D_k}
\Big)
\quad\text{exhaustive}.
\label{eq:learn_cost_sstep}
\end{align}

\paragraph{One sweep.}
Summing \cref{eq:learn_cost_estep,eq:learn_cost_mstep,eq:learn_cost_sstep} over the scales gives
the cost of one outer iteration under the greedy search,
\begin{align}
O\!\left(
\sum_{k=1}^{K}
T_k
\Big[
I_E \big( M_k L + M_k^2 + r_k^3 \big)
+
n_\Theta \big( M_k L + L \bar{C}_k + M_k^2 \big)
+
n_Q (M_k + D_k)\, L D_k \big( \bar{C}_k + B D_k \big)
\Big]
\right).
\label{eq:learn_cost_sweep}
\end{align}
The structure term dominates whenever the candidate set is not screened, which is the setting the
model is usually run in, so candidate screening is the first lever on wall-clock cost. The regime
term $M_k^2$ and the factor term $r_k^3$ matter only at large latent orders.

\paragraph{Memory.}
The dominant arrays are the observed data, the emission tables of shape
$T_k \times M_k \times L$, the pair marginals $\xi_t(m,m')$ of shape
$(T_k - 1) \times M_k \times M_k$, and the smoothed factor covariances of shape
$T_k \times r_k \times r_k$,
\begin{align}
O\!\left(
\sum_{k=1}^{K}
T_k
\big[\, L + M_k L + M_k^2 + r_k^2 \,\big]
\right).
\label{eq:learn_cost_memory}
\end{align}
Beyond the posterior summaries $\bar{\gZ}^{(k+1)}$, which occupy
$O\!\big(T_{k+1}(M_{k+1} + r_{k+1} + L)\big)$, no step needs more than one scale's arrays
resident at a time, so either pass of a sweep can stream the scales.
